% ============================================================================
% BAGIAN AWAL
% ============================================================================

% Halaman judul mengikuti susunan Template Skripsi Informatika UII 2020.
\begin{titlepage}
  \thispagestyle{empty}
  \centering

  {\bfseries\fontsize{16}{24}\selectfont \judulTA\par}

  \vspace{18mm}
  {\bfseries\fontsize{12}{18}\selectfont TUGAS AKHIR\par}

  \vspace{12mm}
  \includegraphics[width=35mm]{UII - Kuning.png}

  \vspace{12mm}
  {\fontsize{12}{18}\selectfont
    Disusun Oleh:\par
    \vspace{5mm}
    \textbf{\namaMahasiswa}\par
    \textbf{\nimMahasiswa}\par
  }

  \vfill
  {\bfseries\fontsize{12}{18}\selectfont
    PROGRAM STUDI INFORMATIKA -- PROGRAM SARJANA\par
    FAKULTAS TEKNOLOGI INDUSTRI\par
    UNIVERSITAS ISLAM INDONESIA\par
    \tahunLulus\par
  }
\end{titlepage}


% Halaman judul dihitung sebagai halaman i tetapi nomornya tidak ditampilkan.
\clearpage
\pagenumbering{roman}
\setcounter{page}{2}
\pagestyle{fancy}

\bagianawaltanpatoc{HALAMAN PENGESAHAN DOSEN PEMBIMBING}

\begin{center}
  {\bfseries\fontsize{16}{24}\selectfont \judulTA\par}
  \vspace{8mm}
  {\bfseries TUGAS AKHIR\par}
  \vspace{6mm}
  Disusun Oleh:\par
  \textbf{\namaMahasiswa}\par
  \textbf{\nimMahasiswa}\par
\end{center}

\vfill
\begin{center}
  \kotaTerbit, \tanggalPengesahan\par
  Dosen Pembimbing,\par
  \vspace{25mm}
  \textbf{\namaPembimbing}\par
  \nipPembimbing
\end{center}

\bagianawaltanpatoc{HALAMAN PENGESAHAN DOSEN PENGUJI}

\begin{center}
  {\bfseries\fontsize{16}{24}\selectfont \judulTA\par}
  \vspace{6mm}
  {\bfseries TUGAS AKHIR\par}
  \vspace{6mm}
  Telah dipertahankan di depan sidang penguji sebagai salah satu syarat untuk memperoleh gelar Sarjana Komputer dari Program Studi Informatika -- Program Sarjana, Fakultas Teknologi Industri, Universitas Islam Indonesia.\par
  \vspace{6mm}
  \kotaTerbit, \tanggalPengesahan
\end{center}

\vspace{8mm}
\noindent
\begin{tabularx}{\textwidth}{@{}Y C@{}}
  \textbf{Tim Penguji} & \textbf{Tanda Tangan} \\
  \addlinespace[6mm]
  \namaPengujiSatu & \rule{40mm}{0.4pt} \\
  \nipPengujiSatu & \\
  \addlinespace[10mm]
  \namaPengujiDua & \rule{40mm}{0.4pt} \\
  \nipPengujiDua & \\
\end{tabularx}

\vfill
\begin{center}
  Mengetahui,\par
  Ketua Program Studi Informatika -- Program Sarjana\par
  Fakultas Teknologi Industri\par
  Universitas Islam Indonesia\par
  \vspace{22mm}
  \textbf{\ketuaProgramStudi}\par
  \nipKetuaProgramStudi
\end{center}

\bagianawaltanpatoc{HALAMAN PERNYATAAN KEASLIAN TUGAS AKHIR}

Yang bertanda tangan di bawah ini:

\vspace{5mm}
\noindent
\begin{tabular}{@{}p{35mm}p{5mm}p{100mm}@{}}
  Nama & : & \namaMahasiswa \\
  NIM & : & \nimMahasiswa \\
\end{tabular}

\vspace{6mm}
Tugas akhir dengan judul:

\begin{center}
  \textbf{\judulTA}
\end{center}

Menyatakan bahwa seluruh komponen dan isi dalam tugas akhir ini adalah hasil karya saya sendiri. Apabila di kemudian hari terbukti bahwa sebagian atau seluruh karya ini bukan hasil karya sendiri, saya bersedia menerima sanksi akademik sesuai dengan ketentuan yang berlaku di Universitas Islam Indonesia.

Demikian surat pernyataan ini dibuat dengan sebenar-benarnya agar dapat dipergunakan sebagaimana mestinya.

\vfill
\begin{flushright}
  \kotaTerbit, \tanggalPengesahan\par
  Yang menyatakan,\par
  \vspace{25mm}
  \textbf{\namaMahasiswa}\par
  \nimMahasiswa
\end{flushright}

\bagianawaltanpatoc{HALAMAN PERSEMBAHAN}
\petunjuk{Halaman ini opsional. Hapus kotak petunjuk dan isi dengan kalimat persembahan yang tidak melanggar etika. Idealnya bagian ini tidak lebih dari satu halaman.}
\begin{center}
  \vfill
  \textit{Tuliskan persembahan di sini.}
  \vfill
\end{center}

\bagianawaltanpatoc{HALAMAN MOTO}
\petunjuk{Halaman ini opsional. Tuliskan moto secara ringkas dan cantumkan sumber apabila mengutip pernyataan orang lain.}
\begin{center}
  \vfill
  \textit{``Tuliskan moto di sini.''}\par
  \vspace{4mm}
  --- Sumber moto
  \vfill
\end{center}

\bagianawal{KATA PENGANTAR}

Puji syukur penulis panjatkan kepada Tuhan Yang Maha Esa karena tugas akhir ini dapat diselesaikan. Tugas akhir ini disusun sebagai salah satu syarat untuk memperoleh gelar Sarjana Komputer pada Program Studi Informatika -- Program Sarjana, Fakultas Teknologi Industri, Universitas Islam Indonesia.

Penulis menyampaikan terima kasih kepada dosen pembimbing, dosen penguji, keluarga, rekan, serta seluruh pihak yang telah memberikan arahan dan dukungan selama penyusunan tugas akhir. Penulis menyadari bahwa dokumen ini masih memiliki keterbatasan sehingga kritik dan saran yang membangun sangat diharapkan.

Semoga tugas akhir ini dapat memberikan manfaat bagi pembaca dan pengembangan ilmu pengetahuan di bidang Informatika.

\vfill
\begin{flushright}
  \kotaTerbit, \bulanLulus\ \tahunLulus\par
  \vspace{18mm}
  \textbf{\namaMahasiswa}
\end{flushright}

\bagianawal{SARI}
\petunjuk{Ganti isi contoh berikut dengan ringkasan yang memuat latar belakang, tujuan atau masalah, metode, hasil utama, dan kesimpulan. Gunakan satu halaman dan tuliskan kata kunci di bagian akhir.}

Paragraf pertama menjelaskan latar belakang dan masalah utama yang melandasi pelaksanaan tugas akhir. Uraian harus singkat, mandiri, dan dapat dipahami tanpa harus membaca keseluruhan laporan.

Paragraf kedua menjelaskan metode, data, perangkat, atau tahapan yang digunakan untuk menyelesaikan masalah. Hindari sitasi, tabel, gambar, dan singkatan yang belum dijelaskan di dalam sari.

Paragraf ketiga menjelaskan hasil utama dan kesimpulan yang diperoleh. Cantumkan hasil kuantitatif penting apabila tersedia dan hindari pernyataan yang tidak didukung oleh hasil tugas akhir.

\vspace{6mm}
\noindent\textbf{Kata kunci:} \kataKunci.

\bagianawal{GLOSARIUM}
\petunjuk{Bagian ini opsional. Urutkan istilah secara alfabetis dan hapus istilah contoh yang tidak digunakan.}

\begin{description}[style=multiline,leftmargin=35mm,font=\normalfont\bfseries]
  \item[Compile] Proses menerjemahkan berkas kode sumber dan berkas terkait menjadi bentuk yang dapat dijalankan atau diproses lebih lanjut.
  \item[Debug] Kegiatan menelusuri dan memperbaiki kesalahan pada perangkat lunak.
  \item[Waterfall] Salah satu model pengembangan perangkat lunak yang membagi proses ke dalam tahapan berurutan.
\end{description}

\clearpage
\phantomsection
\addcontentsline{toc}{chapter}{DAFTAR ISI}
\begingroup
\singlespacing
\tableofcontents
\endgroup

\clearpage
\phantomsection
\addcontentsline{toc}{chapter}{DAFTAR TABEL}
\begingroup
\singlespacing
\listoftables
\endgroup

\clearpage
\phantomsection
\addcontentsline{toc}{chapter}{DAFTAR GAMBAR}
\begingroup
\singlespacing
\listoffigures
\endgroup
